\chapter{Gemelos digitales inalámbricos y trazado de rayos diferenciable}
\label{ch:wdt}

El capítulo anterior estableció la cadena que conduce desde la propagación hasta la
medición de canal. Sobre esa base, este capítulo introduce el gemelo digital como una
realización computacional de dicha cadena. El interés no reside únicamente en reproducir
el canal observado, sino en conservar la relación entre los parámetros físicos de la escena
y las variaciones de la señal que posteriormente se utilizarán para inferir movimiento y
fisiología.

\section{De la simulación de cobertura al gemelo inalámbrico}

Un simulador de propagación predice $H$ a partir de una escena; un gemelo digital persigue, además, correspondencia con un entorno físico concreto, actualización con mediciones y capacidad de responder a cambios. En el contrato de la \cref{eq:thesis_chain}, no basta con generar mapas plausibles: el modelo debe declarar qué escena, frecuencia, antena y observable reproduce. La formulación mínima es
\begin{equation}
\widehat H_m(f,t)=F_m
\left(\Gamma(t),\vect p_T,\vect p_R,f;
\vect s_R,\vect\theta_E\right),
\end{equation}
donde $\theta_E$ incluye parámetros materiales, dispersión, patrones de antena y otros parámetros físicos no conocidos exactamente.

El reto es que una reconstrucción geométrica visualmente correcta puede producir errores de fase considerables. A frecuencia portadora $f_c$, un error de longitud $\delta L$ produce
\begin{equation}
\delta\phi=-2\pi\frac{\delta L}{\lambda}.
\end{equation}
Por tanto, los errores centimétricos o incluso milimétricos pueden ser importantes para el sensado de movimientos pequeños. Esta observación motiva dos líneas complementarias: la calibración diferenciable de propiedades electromagnéticas y la calibración que considera errores locales de fase.

\section{Calibración como problema inverso}

Sea $\vect\theta$ el vector de parámetros del gemelo y $\mathcal D_{\rm cal}$ un conjunto de mediciones alineadas. Se define
\begin{equation}
\vect\theta^\star=\arg\min_{\vect\theta}\;
\sum_{i\in\mathcal D_{\rm cal}}
\calL\left(
H_{m,i}^{\rm real},
F_m(\Gamma_i,\vect p_{T,i},\vect p_{R,i},f;
\vect s_{R,i},\vect\theta)
\right)+\beta R(\vect\theta).
\label{eq:calibration_inverse}
\end{equation}
Para una pérdida real se vectorizan parte real e imaginaria, $\vect u=[\Re H;\Im H]$, y la diferenciabilidad permite obtener
\begin{equation}
\nabla_{\vect\theta}\calL
=\left(\frac{\partial\vect u}{\partial\vect\theta^T}\right)^T
\nabla_{\vect u}\calL.
\end{equation}
Trabajos recientes de calibración con \ac{DRT} muestran que las propiedades materiales, patrones de antena y parámetros de dispersión pueden optimizarse utilizando canales reales \citep{Hoydis2024DiffRT}; otros introducen errores locales de fase para modelar discrepancias geométricas pequeñas que una calibración de potencia no puede absorber \citep{Ruah2024Phase}. En esta tesis se consideran ambas familias como componentes consecutivos, no sustitutos.

\section{Separación de fuentes de error}

Conviene escribir
\begin{equation}
H_{\rm real}-H_{\rm gemelo}=
\Delta H_{\rm geom}+
\Delta H_{\rm mat}+
\Delta H_{\rm scat}+
\Delta H_{\rm ant}+
\Delta H_{\rm hw}+
\epsilon_{\rm model}.
\label{eq:domain_gap_decomp}
\end{equation}
Esta suma es una contabilidad conceptual o una linealización local, no una descomposición identificable a partir de $H$ por sí solo. Un modelo flexible todavía puede usar materiales para compensar coordenadas. Por ello la calibración será jerárquica: registro y sincronización; antenas y escala; materiales y scattering; y sólo al final residuales de fase/modelo. Priors, mediciones auxiliares e intervenciones separadas anclarán cada bloque.

\section{Jerarquía de la literatura y de la evidencia}
\label{sec:literature_evidence_hierarchy}

Los trabajos revisados no se tratarán como soluciones intercambiables. La tesis los
ordena por el operador que modelan, la información que reciben y la evidencia con que
se validan:
\begin{table}[htbp]
\centering\small
\caption{Función científica de la literatura en la cadena metodológica.}
\label{tab:literature_evidence_hierarchy}
\begin{tabularx}{\textwidth}{P{2.6cm}L L}
\toprule
Bloque & Referencias y uso en la tesis & Límite de inferencia\\
\midrule
Calibración estática & Sionna RT, calibración diferenciable y error local de fase
\citep{Hoydis2023SionnaRT,Hoydis2024DiffRT,Ruah2024Phase}. & Mejorar potencia o una
verosimilitud con fase latente no prueba geometría verdadera ni sensado dinámico.\\
Campo y objetos & NeRF$^2$, Learnable WDT, RFScape y RadTwin
\citep{Zhao2023NeRF2,Jiang2025LearnableWDT,Chen2025RFScape,Zhang2026RadTwin}. &
Interpolación de campo, edición de objetos y generalización geométrica son tareas
distintas y requieren entradas equivalentes para compararse.\\
Dinámica simulada & WaveVerse \citep{Zheng2026WaveVerse}. & La coherencia de una
simulación y sus contrastes externos no validan el Jacobiano de nuestro banco ni una
forma de onda respiratoria propia.\\
Inversión & RayLoc, RFDT e InverTwin
\citep{Han2026RayLoc,Chen2026RFDT,Chen2025InverTwin}. & La inversión se evaluará
después de validar el modelo directo, su referencia y sus discontinuidades relevantes.\\
Adquisición y robustez & FarSense, HiSAC, RadioRange y OpenLoRa
\citep{Zeng2019FarSense,Pegoraro2024HiSAC,Lyu2026RadioRange,Mishra2023OpenLoRa}. &
Cocientes, anclas o perturbaciones simuladas dependen de sus supuestos de hardware;
no acreditan coherencia del equipo propio.\\
Representación humana & BreatheSmart, WiMANS, WiMUSE, MURAL-Fi, MM-Fi y WiFi-JEPA
\citep{Mosleh2022BreatheSmart,Huang2024WiMANS,Yi2025WiMUSE,Li2026MURALFI,
Yang2023MMFi,Kim2026WiFiJEPA}. & Una tarea humana bien resuelta no valida por sí
misma la geometría, la fase absoluta ni el gemelo electromagnético.\\
\bottomrule
\end{tabularx}
\end{table}

Esta jerarquía determina el orden de ejecución. DICHASUS, Hoydis y Ruah fijan el
problema del modelo estático; los invariantes y las referencias instrumentales delimitan
el observable; las representaciones de escena se comparan después bajo regímenes de
información comunes; el reflector y el fantoma contrastan fidelidad diferencial; y
la inversión y el aprendizaje humano se activan cuando el observable ha demostrado
sensibilidad. La secuencia no resta valor a los métodos pospuestos: evita usar éxito
en una tarea como evidencia de otra.

\section{DICHASUS dc01 como banco de pruebas del gemelo estático}

La primera validación desarrollada utiliza DICHASUS dc01 y la escena correspondiente empleada en trabajos de calibración diferenciable \citep{Euchner2023DICHASUSdcxx,Hoydis2024DiffRT}. El banco reportado comprende $N=10000$ posiciones y conserva el objeto experimental
\begin{equation}
\mathcal D_{\rm dc01}
=\left\{
\vect p_i,\;H_{m,i}^{\rm real},\;\calP_{m,i}^{\rm RT},\;H_{m,i}^{\rm RT}
\right\}_{i=1}^{N}.
\label{eq:dc01_dataset}
\end{equation}
El soporte geométrico y el conjunto de trayectorias radio--dependiente se conservan como interfaz física, en lugar de almacenar únicamente el canal sintetizado. Esto permite modificar posteriormente materiales, frecuencia, observación y dinámica sin confundir el mundo común $\Gamma$ con una realización particular $\calP_m$.

La comparación B0--B2 y los diagnósticos de fase se documentan en el \cref{ch:avances}. DICHASUS es un banco de sondeo OFDM: su respuesta en frecuencia resulta adecuada para estudiar la propagación, mientras la validación del operador de observación de dispositivos Wi-Fi comerciales requiere pruebas específicas. Conservar la geometría tampoco autoriza reutilizar sin cambios las ganancias y trayectorias efectivas al pasar a LoRa.

\section{Fidelidad de campo y fidelidad diferencial}

Para comunicaciones suele ser suficiente comparar $H_{\rm sim}$ y $H_{\rm real}$. Para el sensado se requiere una condición adicional. Si $q$ es una variable física de interés,
\begin{equation}
J_q=\frac{\partial H}{\partial q},
\end{equation}
y se propone medir
\begin{equation}
E_J=
\frac{\|J_q^{\rm model}-J_q^{\rm ref}\|_F}
{\|J_q^{\rm ref}\|_F+\epsilon}.
\label{eq:jacobian_error}
\end{equation}
Un modelo con bajo \ac{NMSE} de canal, pero con derivadas incorrectas respecto del movimiento humano, puede ser inadecuado para el sensado. Por ello, la tesis distingue entre \emph{fidelidad del modelo directo} y \emph{fidelidad diferencial}.

$J_q^{\rm ref}$ se estimará con perturbaciones mecánicas bidireccionales, paso declarado, repeticiones y alineamiento de fase/ganancia definido antes de observar el resultado. Como el error relativo de la \cref{eq:jacobian_error} es inestable cuando $\|J_q^{\rm ref}\|$ es cercano a cero, se reportará también error absoluto con unidades, correlación compleja y error ponderado por la tarea. El término $\epsilon$ se fijará en las mismas unidades que la norma de referencia y no como constante adimensional arbitraria.

\section{Reconstrucción métrica y registro del gemelo}
\label{sec:geometric_metrology}

La geometría debe representarse con unidades, sistema de coordenadas e incertidumbre
verificables. Se distinguen tres objetivos de diseño: escala global de paredes,
puertas y mobiliario del orden de 1--2 cm; registro de antenas del orden de 2--5 mm;
y referencia local de desplazamiento del orden de 0.1--1 mm. Son rangos de trabajo
que deberán demostrarse mediante controles independientes, no especificaciones
atribuidas a cualquier LiDAR, cámara o actuador. Un experimento con pasos de
0.25 mm sólo es interpretable si su incertidumbre efectiva permite resolverlos.

Se fijará un marco global $W$ y se almacenará, para cada sensor $s$,
\begin{equation}
T_{W\leftarrow s}=\begin{bmatrix}R_{Ws}&\mathbf t_{Ws}\\0&1\end{bmatrix}\in SE(3),
\qquad \widetilde{\mathbf p}_W=T_{W\leftarrow s}\widetilde{\mathbf p}_s,
\label{eq:scene_extrinsics}
\end{equation}
con $R_{Ws}\in SO(3)$. Las transformaciones compuestas deberán respetar esa dirección.
Se registrarán extrínsecos de cámara, LiDAR, radar, transmisores, receptores y antenas;
el punto de montaje no se confundirá con el centro de fase, que puede depender de
frecuencia y dirección. Marcadores rígidos, distancias de control y mediciones repetidas
permitirán verificar el registro. Un residual pequeño de ICP mide concordancia entre
nubes, pero no garantiza exactitud absoluta. Se reservarán, cuando el recinto lo permita,
10--20 distancias de control que no se hayan utilizado para ajustar la reconstrucción.

La comparación propuesta distinguirá cinco condiciones:
\begin{description}
\item[G0.] Planta manual y primitivas geométricas con dimensiones verificadas.
\item[G1.] Reconstrucción RGB-D, con registro y simplificación documentados.
\item[G2.] Reconstrucción LiDAR con las mismas referencias globales.
\item[G3.] Geometría LiDAR regularizada para RF: superficies, espesores, huecos y
objetos relevantes corregidos; materiales nominales y presupuesto RT controlados.
\item[G4.] Geometría anterior con calibración RF de materiales y, si se justifica,
parámetros geométricos restringidos por su incertidumbre metrológica.
\end{description}
G4 cambia más factores que la geometría. Por ello se añadirá una ablación factorial
con materiales nominales/calibrados para separar el efecto del registro, la regularización
y el ajuste electromagnético. Todos los ajustes utilizarán sólo entrenamiento y
validación; los puntos reservados medirán generalización, no servirán para corregir
la escena a posteriori. Se mantendrán mecanismos RT, frecuencias y presupuesto de
muestreo comparables, o se informará explícitamente su coste.

RFDT-Channel ofrece una referencia reciente para construir escenas RF a partir de
RGB y LiDAR y simular canales interiores a 28 GHz \citep{Yao2026RFDTChannel}.
El gemelo de Wang y colaboradores \citep{Wang2026WiFiTwin} contrasta predicción de señal con RSSI real.
Ambos orientan el flujo geométrico, pero una validación de potencia o un canal simulado
no establecen exactitud de fase medida ni fidelidad de micromovimiento de este banco.
La comparación G0--G4 se evaluará sucesivamente en error geométrico, potencia,
retardos, canal relativo y respuesta diferencial. No se exigirá resolver primero
la fase absoluta de toda la habitación para iniciar una prueba local controlada.

Esta distinción cierra el paso entre el modelado electromagnético y la representación: una
aproximación útil debe conservar no sólo valores de canal, sino también las variaciones
que contienen información acerca del estado humano. El capítulo siguiente examina qué
estructuras explícitas, implícitas o híbridas pueden satisfacer ese requisito.
